% =====================================================================
%  CSC_43M04_EP — Computer Vision with Deep Learning (École Polytechnique)
%  Research report: Temporal Action Prediction on SSv2-33 ("What Happens Next?")
%
%  CVPR double-column template.  Requires the official CVPR style files in the
%  same directory (or on TEXINPUTS):  cvpr.sty, cvpr_eso.sty, ieeenat_fullname.bst
%  Download: https://github.com/cvpr-org/author-kit
%
%  Compile:  pdflatex report_cvpr  &&  bibtex report_cvpr  &&  pdflatex x2
%  (or simply  latexmk -pdf report_cvpr.tex)
%
%  Figures live in  docs/figures/  ; this file is expected to sit in  docs/ .
% =====================================================================
\documentclass[10pt,twocolumn,letterpaper]{article}

% ---- CVPR style: use the camera-ready (non-review) mode -------------
\usepackage{cvpr}              % <- final mode (no line numbers, no review banner)
% \usepackage[review]{cvpr}    % <- uncomment for a review version with line numbers

\usepackage{graphicx}
\usepackage{amsmath,amssymb}
\usepackage{booktabs}
\usepackage{xcolor}
\usepackage[pagebackref,breaklinks,colorlinks,allcolors=blue]{hyperref}

\graphicspath{{figures/}}

% ---------------------------------------------------------------------
%  >>> EDIT THESE TWO BLOCKS <<<  (placeholders requested from the user)
% ---------------------------------------------------------------------
\newcommand{\TEAMID}{\textcolor{red}{TEAMID}}            % e.g. 17
\newcommand{\LBONE}{\textcolor{red}{[Stage-1 score / rank]}}  % public-LB top-1 + rank
\newcommand{\LBTWO}{\textcolor{red}{[Stage-2 score / rank]}}  % private-LB top-1 + rank
% Authors: replace the placeholder names below (keep alphabetical order).
% ---------------------------------------------------------------------

\def\cvprPaperID{\TEAMID}
\def\confName{CSC\_43M04\_EP — Modal d'informatique}
\def\confYear{2026}

\title{What Happens Next? Temporal Action Prediction on a 33-Class
Something-Something\,V2 Subset Under a Closed-Track Constraint}

\author{%
Andrea Signoretti\thanks{\textcolor{red}{Replace with the full author list in
alphabetical order; add \texttt{\textbackslash and} between authors. Team ID: \TEAMID.}}\\
École Polytechnique\\
{\tt\small signorettiandrea05@gmail.com}
% ---- example for a second author ----
% \and
% Second Author\\
% École Polytechnique\\
% {\tt\small second@author.fr}
}

\begin{document}
\maketitle

% =====================================================================
\begin{abstract}
We present a systematic study of temporal action prediction on a 33-class
subset of Something-Something\,V2 (SSv2), a benchmark designed so that
classification requires temporal reasoning rather than scene recognition.
Operating under a \emph{closed-track} constraint---training from scratch with
no pretrained weights---and with only four frames per clip, we develop a
progression of architectures from CNN baselines through Temporal Shift Module
(TSM) variants and a hybrid CNN--Transformer. Two findings dominate the
experimental record. First, a 10-epoch \emph{regularisation warm-up} that
defers MixUp/CutMix until basic spatial features have formed yields a
$+36.3$\,pp gain on the Transformer model and is necessary for convergence
across all from-scratch architectures. Second, matching the model's expected
frame count to the true clip length ($T{=}4$ rather than $T{=}5$) gives a
$+4.4$\,pp gain---larger than a $3\times$ parameter scale-up---showing that a
silent data-preparation error can dominate architectural choice in short-clip
video understanding. Combining these insights with focal loss, rotating-fold
training and 10-crop test-time augmentation, a log-softmax late-fusion ensemble
of architecturally complementary models attains \textbf{44.03\% top-1 / 75.37\%
top-5} on the official 6{,}745-clip validation set, $+2.35$\,pp over the best
single model through verified error decorrelation. Failure analysis shows the
residual errors are dominated by real-vs-pretended intent and temporal-direction
confusions that are fundamentally under-constrained at four frames.
\end{abstract}

% =====================================================================
\section{Problem Description}
\label{sec:problem}

\paragraph{Task and dataset.}
The challenge (CSC\_43M04\_EP) is fine-grained action recognition over a
33-class subset of Something-Something\,V2 (SSv2)~\cite{goyal2017ssv2}, framed
as \emph{``What Happens Next?''}. Each clip is a folder of extracted JPEG
frames; models consume a fixed number of frames and output logits over the 33
classes (e.g.\ \emph{Pouring something into something}, \emph{Pretending to
throw something}). The training set contains ${\sim}36{,}000$ clips with a
$20\times$ class imbalance (162--3{,}170 clips per class); the official
validation set (\emph{val-dir}) contains 6{,}745 clips. The evaluation metric
is top-1 accuracy (top-5 reported as a secondary metric).

\paragraph{Why it is hard.}
Three properties make this regime difficult.
\textbf{(i)~Motion-defined labels.} SSv2 was explicitly constructed to defeat
appearance-only recognition~\cite{goyal2017ssv2}: the same objects recur across
classes (any ``something'' can be poured, covered or \emph{pretended} to be
poured), so the discriminative signal lives in the temporal trajectory, not the
static frame---unlike scene-biased benchmarks such as Kinetics~\cite{kay2017kinetics}.
\textbf{(ii)~Extreme temporal sparsity.} Each clip contains exactly
\textbf{four} frames, an order of magnitude fewer than the 8--32 frames standard
in video recognition~\cite{lin2019tsm,bertasius2021timesformer}. Models must
extract maximal temporal information from minimal temporal signal, making the
temporal-aggregation mechanism and the \emph{correctness} of temporal
preprocessing unusually consequential.
\textbf{(iii)~Closed-track constraint.} Pretrained weights (ImageNet, Kinetics)
are prohibited. All spatial features, temporal dynamics and class boundaries
must be learned jointly from ${\sim}36$k clips---roughly $7\times$ fewer than
Kinetics-400---which removes the single largest source of accuracy in modern
video models and amplifies the cost of every design error.

% =====================================================================
\section{Method}
\label{sec:method}

\subsection{Overview and intuition}
Our central hypothesis is that, when pretraining is unavailable, \emph{every}
parameter spent on temporal modelling is a parameter \emph{not} spent on
learning spatial features from scratch. This motivates a preference for
parameter-free or lightweight temporal mechanisms. We therefore evaluate a
ladder of temporal-aggregation strategies on a shared ResNet
backbone~\cite{he2016resnet}, and find that the parameter-free Temporal Shift
Module (TSM)~\cite{lin2019tsm} offers the best accuracy/convergence trade-off,
with a hybrid CNN--Transformer providing a complementary inductive bias for
ensembling. All models share the interface
$(B,T,C,H,W)\!\rightarrow\!(B,\text{33})$.

\subsection{Architectures}
\label{sec:arch}
Figure~\ref{fig:arch} contrasts our two strongest families.

\paragraph{Baselines (own implementation).}
\textbf{CNNBaseline} applies a ResNet18~\cite{he2016resnet} independently to each
of the $T{=}4$ frames and average-pools features over time---a bag-of-frames
lower bound with no temporal reasoning. \textbf{CNNLSTM} replaces the pooling
with a unidirectional LSTM (hidden 256). With only four frames the gating adds
parameters without meaningful gain over pooling, confirming that naive
sequential integration is insufficient.

\paragraph{TSMResNet (our main model; TSM reused).}
We wrap every residual block of a ResNet18 with the Temporal Shift
Module~\cite{lin2019tsm}: a fraction $1/\texttt{fold\_div}$ of channels is
shifted one step backward in time and $1/\texttt{fold\_div}$ forward, letting
frame $t$ read from frames $t{\pm}1$ at \emph{zero} added parameters or FLOPs.
We set $\texttt{fold\_div}{=}4$ (50\% of channels participate per block),
deliberately more aggressive than the default of
8~\cite{lin2019tsm}: with only four frames, context must percolate across a
shallow temporal depth, so a larger shift fraction is required. The shift
operator and the model wrapper are our own implementation following the
published description in~\cite{lin2019tsm}.

\paragraph{VideoFormerLite (VFL; hybrid, our implementation).}
A ResNet18 produces a 512-d feature per frame; a learnable \texttt{[CLS]}
token~\cite{dosovitskiy2021vit} and temporal positional embeddings are
prepended, and the sequence passes through a 2-layer \emph{pre-LN} Transformer
encoder ($d{=}512$, 8 heads)~\cite{xiong2020prenorm}. The \texttt{[CLS]} output
is projected to logits. With $T{=}4$ tokens the attention block is
computationally trivial (${\sim}0.03$\,GFLOPs vs.\ ${\sim}7.2$ for the backbone);
it acts as a \emph{global} temporal aggregator---any frame pair is one attention
step apart---in contrast to TSM's \emph{local} shift. This is a deliberately
lightweight relative of TimeSformer~\cite{bertasius2021timesformer}: we keep a
CNN spatial encoder (more sample-efficient from scratch) and use the Transformer
only as a temporal head.

\paragraph{Negative-result architectures.}
We also trained \textbf{R(2+1)D-18}~\cite{tran2018r2plus1d} (torchvision,
from scratch) and a \textbf{TSM-ResNet50 + BiGRU} variant. Both are reported in
Section~\ref{sec:results} as informative failures.

\begin{figure}[t]
  \centering
  \includegraphics[width=\linewidth]{arch_comparison.png}
  \caption{The two complementary temporal mechanisms ensembled in this work.
  \textbf{TSM} (left) mixes time \emph{locally} via parameter-free channel shifts
  inside every ResNet block; \textbf{VideoFormerLite} (right) mixes time
  \emph{globally} via \texttt{[CLS]}-token self-attention over the four frame
  tokens. Their orthogonal inductive biases drive the ensemble gain
  (Sec.~\ref{sec:ablation}).}
  \label{fig:arch}
\end{figure}

\subsection{Training recipe}
\label{sec:training}
All models are trained from scratch with mixed-precision
(\texttt{torch.autocast}) on a single NVIDIA RTX A4000 (16\,GB), managed with
Hydra configs. The recipe (Table~\ref{tab:recipe}) reflects two
closed-track-specific findings.

\paragraph{Regularisation warm-up.}
MixUp~\cite{zhang2018mixup} and CutMix~\cite{yun2019cutmix} assume a network that
already produces meaningful class probabilities. Applied from epoch 1 to a
randomly initialised model they create conflicting gradients on top of basic
feature learning, preventing convergence. We therefore \emph{defer} all
label-mixing for the first 10 epochs (\emph{regularisation warm-up}). This single
change is the largest controlled effect we measured (Sec.~\ref{sec:ablation}).

\paragraph{Frame-count correctness.}
The dataset's uniform \texttt{linspace} sampler returns duplicated indices when
the requested $T$ exceeds the four available frames. Setting $T{=}4$ ensures
every sampled frame is unique; duplicated tokens otherwise carry no new
information yet still corrupt the temporal-shift signal (Sec.~\ref{sec:ablation}).

\paragraph{Loss and optimisation.}
The best model uses focal loss~\cite{lin2017focal} ($\gamma{=}2$) with label
smoothing $0.1$: SSv2 deliberately groups near-identical
classes~\cite{goyal2017ssv2}, and the $(1{-}p_t)^\gamma$ reweighting
concentrates gradient on these hard within-group pairs. We optimise with
AdamW~\cite{loshchilov2019adamw} (peak lr $10^{-3}$, weight decay $5\times10^{-4}$),
an 8-epoch linear warm-up, and cosine annealing~\cite{loshchilov2017sgdr}.
Horizontal flipping is \emph{disabled} because it inverts the
left/right-directional SSv2 labels; spatial augmentation
(RandAugment~\cite{cubuk2020randaugment}, random erasing~\cite{zhong2020randomerase})
is applied \emph{temporally consistently} (identical parameters across the four
frames) to preserve motion. Class imbalance is handled by
$\sqrt{\,}$-frequency weighted sampling. Stochastic depth~\cite{huang2016stochdepth}
(rate 0.1) and head dropout (0.3) regularise the network.

\paragraph{Rotating $k$-fold and inference.}
For the strongest variants we train with rotating 5-fold splits (each epoch sees
80\% of the data, epochs scaled by $1/0.8$), which both regularises and yields
decorrelated checkpoints for ensembling. At inference we optionally apply
10-crop test-time augmentation (TTA: 5 spatial crops $\times$ 2 flips) and
average logits. Ensembling combines per-model TTA-averaged log-probabilities
(Sec.~\ref{sec:ensemble}).

\paragraph{Reused vs.\ own code.}
\emph{Reused:} the ResNet18/R(2+1)D backbones (torchvision), the TSM
\emph{concept}~\cite{lin2019tsm}, and standard augmentation/optimiser primitives.
\emph{Our own:} the TSM block wrapper, VideoFormerLite, the full Hydra training
pipeline (\texttt{train.py}/\texttt{train\_kfold.py}), the temporally-consistent
augmentation, the dataset/sampler, and all evaluation and ensembling code.

\begin{table}[t]
\centering\small
\setlength{\tabcolsep}{4pt}
\begin{tabular}{@{}lll@{}}
\toprule
Technique & Setting & Source \\
\midrule
Reg.\ warm-up      & 10 ep, no MixUp/CutMix & \textbf{this work} \\
MixUp              & $\alpha{=}0.2$, after ep 10 & \cite{zhang2018mixup} \\
CutMix             & $p{=}0.25$, after ep 10 & \cite{yun2019cutmix} \\
Focal loss         & $\gamma{=}2$ & \cite{lin2017focal} \\
Label smoothing    & $\varepsilon{=}0.1$ & \cite{muller2019labelsmooth} \\
RandAugment        & 2 ops, mag.\ 0.5 & \cite{cubuk2020randaugment} \\
Random erasing     & $p{=}0.25$ & \cite{zhong2020randomerase} \\
Weighted sampling  & power $0.5$ & --- \\
Horizontal flip    & \textbf{disabled} & label semantics \\
Stochastic depth   & rate $0.1$ & \cite{huang2016stochdepth} \\
Optimiser          & AdamW, lr $10^{-3}$, wd $5{\times}10^{-4}$ & \cite{loshchilov2019adamw} \\
Schedule           & 8-ep warm-up + cosine & \cite{loshchilov2017sgdr} \\
\bottomrule
\end{tabular}
\caption{Training recipe for the best single model (TSM-Ultra-v2, ResNet18,
$T{=}4$, \texttt{fold\_div}=4, 150 epochs, ${\sim}5.8$\,h on one A4000).}
\label{tab:recipe}
\end{table}

% =====================================================================
\section{Results}
\label{sec:results}

\subsection{Experimental protocol}
We report top-1 accuracy on the official val-dir set (6{,}745 clips), evaluated
without TTA unless stated, to isolate model quality from inference-time
augmentation. An internal 20\% split of the training data is used for model
selection; it runs $15$--$19$\,pp higher than val-dir because of distribution
shift, so val-dir is treated as ground truth. The seed is fixed (42).

\paragraph{Leaderboard.}
On the official challenge leaderboard our submission scored \LBONE{} in Stage~1
and \LBTWO{} in Stage~2. These correspond to the log-softmax ensemble described
in Sec.~\ref{sec:ensemble}; val-dir numbers below are the offline counterpart.

\subsection{Main results}
Table~\ref{tab:main} and Figure~\ref{fig:permodel} summarise the progression.
The best \emph{single} model, \textbf{TSM-Ultra-v2}, reaches \textbf{38.25\%}
val-dir top-1 (57.73\% internal); the best \emph{ensemble} reaches
\textbf{44.03\%} top-1 / 75.37\% top-5.

\begin{table*}[t]
\centering\small
\setlength{\tabcolsep}{5pt}
\begin{tabular}{@{}lllccccc@{}}
\toprule
Model & Backbone & Temporal module & $T$ & Params & Int.\ Val & Val-Dir T1 & Val-Dir T5 \\
\midrule
CNNBaseline             & ResNet18 & Avg pool            & 4 & 11.2M & ---     & ---     & --- \\
CNNLSTM                 & ResNet18 & LSTM ($h{=}256$)    & 4 & 11.5M & marg.   & ---     & --- \\
R(2+1)D-18 (ep15)       & R(2+1)D  & 3D-factored conv    & 8 & 31M   & 10.78\% & ---     & --- \\
VFL (no warm-up)        & ResNet18 & 2-layer Transformer & 4 & 17.5M & \phantom{0}9.45\% & --- & --- \\
VFL (warm-up, 60ep)     & ResNet18 & 2-layer Transformer & 4 & 17.5M & 45.78\% & ---     & --- \\
TSM-Ultra               & ResNet18 & Channel shift       & 5 & 11.2M & 53.32\% & 34.91\% & 66.88\% \\
TSM-Ultra-50            & ResNet50 & Channel shift       & 8 & 23.6M & 45.26\% & ---     & --- \\
VFL-Ultra               & ResNet18 & 2-layer Transformer & 4 & 17.5M & 53.37\% & 32.11\% & 64.85\% \\
\textbf{TSM-Ultra-v2}   & ResNet18 & Channel shift + focal & 4 & 11.2M & \textbf{57.73\%} & \textbf{38.25\%} & \textbf{68.70\%} \\
\midrule
3-model weighted ens.\ (+TTA)  & --- & Weighted logit avg & --- & --- & --- & 40.43\% & 71.27\% \\
\textbf{Best ensemble} (+rot+TTA) & --- & Log-softmax avg & --- & --- & --- & \textbf{44.03\%} & \textbf{75.37\%} \\
\bottomrule
\end{tabular}
\caption{Main results. Internal-val is a 20\% held-out split (optimistically
biased); val-dir is the official 6{,}745-clip set. The best single model fixes
the frame count and adds focal loss; the best ensemble fuses TSM and VFL with
rotating folds and 10-crop TTA.}
\label{tab:main}
\end{table*}

\begin{figure}[t]
  \centering
  \includegraphics[width=\linewidth]{fig01_per_model_accuracy.png}
  \caption{Per-model val-dir top-1/top-5. The TSM family dominates; correcting
  the frame count (TSM-Ultra$\rightarrow$TSM-Ultra-v2) is the single largest
  per-model jump.}
  \label{fig:permodel}
\end{figure}

\subsection{Ablations}
\label{sec:ablation}

\paragraph{Regularisation warm-up ($+36.3$\,pp).}
Table~\ref{tab:warmup} isolates the warm-up: identical architecture and config,
single variable. Without it, VFL stalls at 9.45\% (killed at epoch 4) and
R(2+1)D at 10.78\%; with a 10-epoch warm-up VFL reaches 45.78\%. The effect
transfers across architectures, suggesting it is a general property of
cold-start training with label-mixing augmentations rather than an
architecture-specific quirk.

\begin{table}[t]
\centering\small
\begin{tabular}{@{}lccc@{}}
\toprule
Run & Warm-up & Val @ep10 & Final \\
\midrule
VFL (killed)        & none      & ${\sim}9.4\%$ & \phantom{0}9.45\% \\
VFL                 & 10 ep     & 17.58\%       & \textbf{45.78\%} (60ep) \\
TSM-Ultra           & 10 ep     & 25.60\%       & \textbf{53.32\%} (80ep) \\
R(2+1)D (killed)    & none      & ${\sim}8.4\%$ & 10.78\% (15ep) \\
\bottomrule
\end{tabular}
\caption{Regularisation-warm-up ablation. Deferring MixUp/CutMix by 10 epochs is
the largest controlled effect in this study ($+36.3$\,pp on VFL).}
\label{tab:warmup}
\end{table}

\paragraph{Frame count ($+4.4$\,pp).}
Reducing $T$ from 5 to 4 to match the true clip length improves internal val by
$+4.41$\,pp ($53.32\!\rightarrow\!57.73$) and val-dir by $+3.34$\,pp
($34.91\!\rightarrow\!38.25$). Crucially, scaling instead to ResNet50 \emph{with}
$T{=}8$ (TSM-Ultra-50) \emph{regresses} to 45.26\%: a $3\times$ parameter
increase cannot compensate for corrupted temporal input
(Fig.~\ref{fig:permodel}). The correct order is \emph{fix the data, then scale}.

\paragraph{Rotating folds and TTA.}
Rotating 5-fold training adds ${\sim}+1.2$\,pp to each single model
(e.g.\ TSM-v2 $38.25\!\rightarrow\!39.45$); 10-crop TTA adds a further
${\sim}+2.6$\,pp (TSM-v2 $\rightarrow\!40.86$, TSM-v2-rot $\rightarrow\!41.68$),
the strongest single configuration before ensembling (Fig.~\ref{fig:ttarot}).

\begin{figure}[t]
  \centering
  \includegraphics[width=\linewidth]{fig02_tta_effect.png}
  \caption{Effect of 10-crop TTA across models: a consistent
  ${\sim}+2.5$--$3$\,pp top-1 gain, larger on top-5.}
  \label{fig:ttarot}
\end{figure}

\paragraph{Ensemble components (leave-one-out).}
On the three-model weighted sub-ensemble (40.43\%), removing any component
degrades accuracy (Table~\ref{tab:loo}, Fig.~\ref{fig:loo}): no member is
redundant. The largest drop ($-3.17$\,pp) comes from removing the strongest model
(TSM-Ultra-v2), but removing VFL---the \emph{weakest} member---still costs
$-0.61$\,pp, confirming that its globally-attentive errors decorrelate from
TSM's locally-shifted ones.

\begin{table}[t]
\centering\small
\begin{tabular}{@{}lcc@{}}
\toprule
Configuration & Top-1 & $\Delta$ \\
\midrule
Full ensemble (A+B+C)    & \textbf{40.43\%} & --- \\
\;$-$ TSM-Ultra-v2 (A)   & 37.26\% & $-3.17$ \\
\;$-$ TSM-Ultra (B)      & 39.56\% & $-0.87$ \\
\;$-$ VFL-Ultra (C)      & 39.82\% & $-0.61$ \\
\bottomrule
\end{tabular}
\caption{Leave-one-out on the weighted sub-ensemble. Every component contributes
positively---ensemble gain stems from error decorrelation, not accuracy
averaging.}
\label{tab:loo}
\end{table}

\begin{figure}[t]
  \centering
  \includegraphics[width=\linewidth]{fig11_leave_one_out.png}
  \caption{Leave-one-out ensemble study (val-dir top-1). All three members are
  necessary; the weakest standalone model (VFL) still adds value through
  decorrelation.}
  \label{fig:loo}
\end{figure}

\subsection{The final ensemble}
\label{sec:ensemble}
The best closed-track result fuses \texttt{tsm\_v2}, its rotating-fold variant
\texttt{tsm\_v2\_rot}, and \texttt{vfl\_rot}, each with 10-crop TTA, by averaging
\emph{log-softmax} probabilities (Fig.~\ref{fig:ensemble}). Log-space averaging
down-weights over-confident wrong predictions and rewards agreement, which suits
members with different calibration (focal vs.\ CE; single-split vs.\ rotating
fold). It reaches \textbf{44.03\% / 75.37\%}, $+2.35$\,pp over the best single
configuration (41.68\%). Adding the weaker closed models (plain \texttt{tsm},
\texttt{vfl}) \emph{hurts} ($\rightarrow\!42.73\%$): their noisier signals dilute
the two strong TSM predictions, a concrete instance of the
``diversity-not-quantity'' principle.

\begin{figure}[t]
  \centering
  \includegraphics[width=\linewidth]{fig04_ensemble.png}
  \caption{Ensemble vs.\ its components. Log-softmax late fusion of two TSM
  variants and one VFL variant (all +TTA) exceeds the best single model by
  $+2.35$\,pp top-1.}
  \label{fig:ensemble}
\end{figure}

\subsection{Per-class and failure analysis}
\label{sec:failure}
The ensemble's macro-F1 is 0.358, but per-class F1 spans a wide range
(Fig.~\ref{fig:perclass}). The \textbf{best} classes (F1\,$\geq$\,0.54) are all
\emph{clean directional single-object motions}---\emph{Pulling left$\rightarrow$right}
(0.62), \emph{Pouring into} (0.58), \emph{Moving closer} (0.57),
\emph{Uncovering} (0.55), \emph{Folding} (0.53)---precisely the patterns TSM's
per-block shift is built to encode.

The \textbf{worst} classes expose two systematic failure modes
(Figs.~\ref{fig:confmat},~\ref{fig:pairs}). \emph{(i)~Intent/``pretending''
ambiguity}: \emph{Picking something up} (F1\,$=$\,0.07) is confused 30/199 times
with \emph{Pretending to pick something up}---the executed motion is visually
identical and only the (often off-screen) outcome differs. \emph{(ii)~Temporal
direction \& near-duplicate pairs}: \emph{Folding}$\leftrightarrow$\emph{Unfolding}
(29 and 24 mutual confusions) and \emph{Moving up}$\rightarrow$\emph{Pretending to
pick up} (40) differ only in motion direction or completion, which four frames
under-determine. The lowest score, \emph{Spilling something next to something}
(F1\,$=$\,0.03, $n{=}60$), combines visual ambiguity with low support. These
residual errors coincide exactly with SSv2's design intent~\cite{goyal2017ssv2}:
they are dataset-level limits of a four-frame appearance model, not tuning
failures.

\begin{figure}[t]
  \centering
  \includegraphics[width=\linewidth]{fig05_per_class_f1.png}
  \caption{Per-class F1 (ensemble, val-dir). Directional single-object motions
  score highest; ``pretending'' and direction-ambiguous classes score lowest.}
  \label{fig:perclass}
\end{figure}

\begin{figure}[t]
  \centering
  \includegraphics[width=\linewidth]{fig06_confusion_matrix.png}
  \caption{Confusion matrix (ensemble, val-dir). Off-diagonal mass concentrates
  on semantically adjacent pairs (real vs.\ pretended; fold vs.\ unfold).}
  \label{fig:confmat}
\end{figure}

\begin{figure}[t]
  \centering
  \includegraphics[width=\linewidth]{fig07_confused_pairs.png}
  \caption{Most-confused class pairs. The dominant errors are intent
  (real vs.\ pretended) and temporal-direction confusions.}
  \label{fig:pairs}
\end{figure}

\subsection{What did not work}
\textbf{Scaling before fixing data} (TSM-Ultra-50): $3\times$ parameters and
$T{=}8$ regressed 8\,pp below TSM-Ultra-v2 because $T{=}8$ re-introduced
duplicated frames. \textbf{R(2+1)D from scratch}: 3D-factored convolutions need
far more data/pretraining; combined with a forced resolution drop (112$\times$112
for VRAM) it stalled at 10.78\%. \textbf{TSM-ResNet50 + BiGRU}: a heavy backbone,
a sequential head and a rotating-fold split (which cuts effective per-epoch data
$5\times$) together failed to converge (${\sim}7.6\%$). The common lesson is that
under a closed track, convergence---not capacity---is the binding constraint.

\subsection{Comparison to published SSv2 baselines}
Our best single model (38.25\%) sits ${\sim}20.9$\,pp below the lowest published
TSM result (ResNet50, Kinetics-pretrained, $T{=}8$: 59.1\%~\cite{lin2019tsm}).
We attribute this gap to no pretraining (${\sim}15$\,pp), a smaller backbone
(${\sim}3$\,pp) and fewer frames (${\sim}2$\,pp)---i.e.\ the model performs as
first-principles reasoning predicts given the constraints, and the 44.03\%
ensemble narrows the gap further. State-of-the-art open-track numbers
(Video Swin~\cite{liu2022videoswin} 69.6\%; VideoMAE\,V2~\cite{wang2023videomae2}
74.7\% on the full 174-class set) quantify the value of large-scale pretraining
that the closed track forbids.

% =====================================================================
\section{Conclusion and Future Work}
\label{sec:conclusion}
We documented a systematic, closed-track progression from a CNN baseline to a
log-softmax ensemble for four-frame action prediction on SSv2-33. Two
data-centric findings dominate: a 10-epoch regularisation warm-up that unlocks
convergence for from-scratch label-mixing ($+36.3$\,pp), and a frame-count
correction ($T{=}5\!\rightarrow\!4$) that beats a $3\times$ parameter scale-up
($+4.4$\,pp). Built on these, a log-softmax fusion of two TSM variants and a
VideoFormerLite---chosen for orthogonal local-shift vs.\ global-attention
biases, with rotating folds and 10-crop TTA---achieves \textbf{44.03\% top-1 /
75.37\% top-5} on val-dir, with a leave-one-out study confirming every component
contributes through error decorrelation. Failure analysis shows the residual
errors are exactly the real-vs-pretended and temporal-direction confusions that
SSv2 was designed to demand and that four frames under-determine.

\paragraph{Future work.}
(i)~Isolate the focal-loss contribution, currently confounded with the
frame-count fix. (ii)~Multi-clip temporal inference (3--5 windows/video,
${\sim}+1$--$2$\,pp expected). (iii)~Frame interpolation ($4\!\rightarrow\!8$) as
synthetic temporal augmentation, e.g.\ with a MaxViT~\cite{tu2022maxvit} head.
(iv)~An explicit open- vs.\ closed-track study to quantify the pretraining gap on
this subset. (v)~Stratified cross-validation for more robust model selection
under the $20\times$ class imbalance.

% =====================================================================
\section*{Contributions}
\noindent\textcolor{red}{\emph{[Fill in per team member. Template:]}}
\textbf{Andrea Signoretti} --- designed and implemented the TSM and
VideoFormerLite models, the Hydra training/k-fold pipeline, the ensemble and TTA
code, ran all experiments and ablations, and produced the analysis figures and
this report. \emph{(For a multi-person team, split: model implementation,
training-pipeline engineering, ensemble/evaluation, data analysis \& figures,
and writing, attributing each explicitly.)}

% =====================================================================
{\small
\bibliographystyle{ieeenat_fullname}
\bibliography{report_cvpr}
}

\end{document}
